\documentclass[10pt,a4paper]{amsart}
\usepackage{amsmath,amssymb,amsthm,graphicx}
\usepackage{verbatim}
\newenvironment{reference}{\comment}{\endcomment}
%\newenvironment{reference}{}{}
\newenvironment{slogan}{\comment}{\endcomment}
\newenvironment{history}{\comment}{\endcomment}

% For commutative diagrams we use Xy-pic
\usepackage[all]{xy}

% We use 2cell for 2-commutative diagrams.
\xyoption{2cell}
% Theorem environments.
%
\theoremstyle{plain}
\newtheorem{theorem}[subsection]{Theorem}
\newtheorem{proposition}[subsection]{Proposition}
\newtheorem{lemma}[subsection]{Lemma}

\theoremstyle{definition}
\newtheorem{definition}[subsection]{Definition}
\newtheorem{example}[subsection]{Example}
\newtheorem{exercise}[subsection]{Exercise}
\newtheorem{situation}[subsection]{Situation}

\theoremstyle{remark}
\newtheorem{remark}[subsection]{Remark}
\newtheorem{remarks}[subsection]{Remarks}

\numberwithin{equation}{subsection}

% Macros
%
\def\lim{\mathop{\mathrm{lim}}\nolimits}
\def\colim{\mathop{\mathrm{colim}}\nolimits}
\def\Spec{\mathop{\mathrm{Spec}}}
\def\Hom{\mathop{\mathrm{Hom}}\nolimits}
\def\Ext{\mathop{\mathrm{Ext}}\nolimits}
\def\SheafHom{\mathop{\mathcal{H}\!\mathit{om}}\nolimits}
\def\SheafExt{\mathop{\mathcal{E}\!\mathit{xt}}\nolimits}
\def\Sch{\mathit{Sch}}
\def\Mor{\mathop{\mathrm{Mor}}\nolimits}
\def\Ob{\mathop{\mathrm{Ob}}\nolimits}
\def\Sh{\mathop{\mathit{Sh}}\nolimits}
\def\NL{\mathop{N\!L}\nolimits}
\def\CH{\mathop{\mathrm{CH}}\nolimits}
\def\proetale{{pro\text{-}\acute{e}tale}}
\def\etale{{\acute{e}tale}}
\def\QCoh{\mathit{QCoh}}
\def\Ker{\mathop{\mathrm{Ker}}}
\def\Im{\mathop{\mathrm{Im}}}
\def\Coker{\mathop{\mathrm{Coker}}}
\def\Coim{\mathop{\mathrm{Coim}}}

% Boxtimes
%
\DeclareMathSymbol{\boxtimes}{\mathbin}{AMSa}{"02}

%
% Macros for moduli stacks/spaces
%
\def\QCohstack{\mathcal{QC}\!\mathit{oh}}
\def\Cohstack{\mathcal{C}\!\mathit{oh}}
\def\Spacesstack{\mathcal{S}\!\mathit{paces}}
\def\Quotfunctor{\mathrm{Quot}}
\def\Hilbfunctor{\mathrm{Hilb}}
\def\Curvesstack{\mathcal{C}\!\mathit{urves}}
\def\Polarizedstack{\mathcal{P}\!\mathit{olarized}}
\def\Complexesstack{\mathcal{C}\!\mathit{omplexes}}
% \Pic is the operator that assigns to X its picard group, usage \Pic(X)
% \Picardstack_{X/B} denotes the Picard stack of X over B
% \Picardfunctor_{X/B} denotes the Picard functor of X over B
\def\Pic{\mathop{\mathrm{Pic}}\nolimits}
\def\Picardstack{\mathcal{P}\!\mathit{ic}}
\def\Picardfunctor{\mathrm{Pic}}
\def\Deformationcategory{\mathcal{D}\!\mathit{ef}}
\usepackage{cleveref}
\crefname{originallemma}{original lemma}{original lemmas}
\Crefname{originallemma}{Original lemma}{Original lemmas}
\makeatletter
\newlabel{algebra-lemma-completion-generalities}{{0315}{source}{}{sourcepointer-1}{}}
\newlabel{algebra-lemma-completion-generalities@cref}{{[originallemma][0][]0315}{[1][0][]source}{}{}{}}
\makeatother
\crefname{originalequation}{original equation}{original equations}
\Crefname{originalequation}{Original equation}{Original equations}
\makeatletter
\newlabel{equation-C-prime}{{0AL4}{source}{}{sourcepointer-2}{}}
\newlabel{equation-C-prime@cref}{{[originalequation][0][]0AL4}{[1][0][]source}{}{}{}}
\makeatother
\crefname{originalreference}{original reference}{original references}
\Crefname{originalreference}{Original reference}{Original references}
\makeatletter
\newlabel{lemma-NL-is-completion}{{0GAF}{source}{}{sourcepointer-3}{}}
\newlabel{lemma-NL-is-completion@cref}{{[originalreference][0][]0GAF}{[1][0][]source}{}{}{}}
\makeatother
\crefname{originalreference}{original reference}{original references}
\Crefname{originalreference}{Original reference}{Original references}
\makeatletter
\newlabel{lemma-equivalent-with-artin-smooth}{{0GAJ}{source}{}{sourcepointer-4}{}}
\newlabel{lemma-equivalent-with-artin-smooth@cref}{{[originalreference][0][]0GAJ}{[1][0][]source}{}{}{}}
\makeatother
\crefname{originalreference}{original reference}{original references}
\Crefname{originalreference}{Original reference}{Original references}
\makeatletter
\newlabel{lemma-exact-sequence-NL}{{00S2}{source}{}{sourcepointer-5}{}}
\newlabel{lemma-exact-sequence-NL@cref}{{[originalreference][0][]00S2}{[1][0][]source}{}{}{}}
\makeatother
\crefname{originalreference}{original reference}{original references}
\Crefname{originalreference}{Original reference}{Original references}
\makeatletter
\newlabel{lemma-get-morphism-general-better}{{0GAQ}{source}{}{sourcepointer-6}{}}
\newlabel{lemma-get-morphism-general-better@cref}{{[originalreference][0][]0GAQ}{[1][0][]source}{}{}{}}
\makeatother
\crefname{originalreference}{original reference}{original references}
\Crefname{originalreference}{Original reference}{Original references}
\makeatletter
\newlabel{lemma-presentation-rig-smooth}{{0GAV}{source}{}{sourcepointer-7}{}}
\newlabel{lemma-presentation-rig-smooth@cref}{{[originalreference][0][]0GAV}{[1][0][]source}{}{}{}}
\makeatother
\crefname{originalsection}{original section}{original sections}
\Crefname{originalsection}{Original section}{Original sections}
\makeatletter
\newlabel{local-cohomology-section-uniform}{{0GA6}{source}{}{sourcepointer-8}{}}
\newlabel{local-cohomology-section-uniform@cref}{{[originalsection][0][]0GA6}{[1][0][]source}{}{}{}}
\makeatother
\crefname{originallemma}{original lemma}{original lemmas}
\Crefname{originallemma}{Original lemma}{Original lemmas}
\makeatletter
\newlabel{more-algebra-lemma-base-change-property-ext-1-annihilated}{{0G9H}{source}{}{sourcepointer-9}{}}
\newlabel{more-algebra-lemma-base-change-property-ext-1-annihilated@cref}{{[originallemma][0][]0G9H}{[1][0][]source}{}{}{}}
\makeatother
\crefname{originallemma}{original lemma}{original lemmas}
\Crefname{originallemma}{Original lemma}{Original lemmas}
\makeatletter
\newlabel{more-algebra-lemma-quotient-by-idempotent}{{0ALR}{source}{}{sourcepointer-10}{}}
\newlabel{more-algebra-lemma-quotient-by-idempotent@cref}{{[originallemma][0][]0ALR}{[1][0][]source}{}{}{}}
\makeatother
\crefname{originallemma}{original lemma}{original lemmas}
\Crefname{originallemma}{Original lemma}{Original lemmas}
\makeatletter
\newlabel{more-algebra-lemma-two-term-base-change}{{0G9G}{source}{}{sourcepointer-11}{}}
\newlabel{more-algebra-lemma-two-term-base-change@cref}{{[originallemma][0][]0G9G}{[1][0][]source}{}{}{}}
\makeatother
\title{An approximate presentation with divisorial Jacobian bound algebraizes to the completion of a finite type algebra}
\author{The Stacks Project Authors}
\date{stacks-project:restricted.tex:local-snapshot-20261012; GNU FDL 1.2 or later}
\begin{document}
\maketitle
\noindent\textit{Editorial note.} All delivered mathematical text is the author's own, copied from the pinned Stacks Project chapter file restricted.tex; no statement or proof step was rewritten, reordered or removed.\par\medskip
\noindent\textit{Editorial note.} The counted proof is the author's complete argument: a rig-smooth reduction of the hypotheses, an engineered finite type algebra whose Jacobian determinant gives a truncation-independent annihilation bound, a lift of the mod-I\textasciicircum{}n presentation to a genuine surjection, and a kernel control finished by a free-module rank computation after inverting b\par\medskip
\noindent\textit{Editorial note.} The label lemma-approximate-presentation-rig-smooth names the counted statement, and the proof environment immediately following it in the source is delivered verbatim as its proof.\par\medskip
\noindent\textit{Editorial note.} The counted statement is carried by the author's own lemma environment in the chapter source.\par\medskip
\noindent\textit{Allowed standard prerequisites (editorial).}\begin{itemize}
\item The category of topological A-algebras complete along I for a Noetherian ring A with ideal I (the displayed category in the source section), and the rig-smooth property over the pair (A, I)
\item The completed naive cotangent complex NL and its long exact sequence, as developed in the same chapter
\item The uniformity integers d and q of the associated graded Gr\_I(B) from the local cohomology section cited in the proof
\end{itemize}\medskip
\ifdefined\hypertarget\hypertarget{sourcepointer-1}{}\fi
\noindent\textit{Source reference.} Original lemma 0315: sources/stacks-chow/tags.txt, source lines 3972--3972. Prerequisite pointer invoked at proof line(s) 1726; it names a standard fact of a sibling chapter that the counted argument uses, and every step of the counted proof built on it is retained verbatim in the delivered body.\par\medskip
\ifdefined\hypertarget\hypertarget{sourcepointer-2}{}\fi
\noindent\textit{Source reference.} Original equation 0AL4: sources/stacks-restricted/pinned-original-source.tex, source lines 326--326. Prerequisite pointer invoked at proof line(s) 1640; it names a standard fact of the same chapter that the counted argument uses, and every step of the counted proof built on it is retained verbatim in the delivered body.\par\medskip
\ifdefined\hypertarget\hypertarget{sourcepointer-3}{}\fi
\noindent\textit{Source reference.} Original reference 0GAF: sources/stacks-restricted/pinned-original-source.tex, source lines 560--560. Prerequisite pointer invoked at proof line(s) 1697; it names a standard fact of the same chapter that the counted argument uses, and every step of the counted proof built on it is retained verbatim in the delivered body.\par\medskip
\ifdefined\hypertarget\hypertarget{sourcepointer-4}{}\fi
\noindent\textit{Source reference.} Original reference 0GAJ: sources/stacks-restricted/pinned-original-source.tex, source lines 843--843. Prerequisite pointer invoked at proof line(s) 1657; it names a standard fact of the same chapter that the counted argument uses, and every step of the counted proof built on it is retained verbatim in the delivered body.\par\medskip
\ifdefined\hypertarget\hypertarget{sourcepointer-5}{}\fi
\noindent\textit{Source reference.} Original reference 00S2: sources/stacks-restricted/pinned-original-source.tex, source lines 725--725. Prerequisite pointer invoked at proof line(s) 1744; it names a standard fact of the same chapter that the counted argument uses, and every step of the counted proof built on it is retained verbatim in the delivered body.\par\medskip
\ifdefined\hypertarget\hypertarget{sourcepointer-6}{}\fi
\noindent\textit{Source reference.} Original reference 0GAQ: sources/stacks-restricted/pinned-original-source.tex, source lines 1167--1167. Prerequisite pointer invoked at proof line(s) 1718; it names a standard fact of the same chapter that the counted argument uses, and every step of the counted proof built on it is retained verbatim in the delivered body.\par\medskip
\ifdefined\hypertarget\hypertarget{sourcepointer-7}{}\fi
\noindent\textit{Source reference.} Original reference 0GAV: sources/stacks-restricted/pinned-original-source.tex, source lines 1598--1598. Prerequisite pointer invoked at proof line(s) 1674; it names a standard fact of the same chapter that the counted argument uses, and every step of the counted proof built on it is retained verbatim in the delivered body.\par\medskip
\ifdefined\hypertarget\hypertarget{sourcepointer-8}{}\fi
\noindent\textit{Source reference.} Original section 0GA6: sources/stacks-chow/tags.txt, source lines 20326--20326. Prerequisite pointer invoked at proof line(s) 1717; it names a standard fact of a sibling chapter that the counted argument uses, and every step of the counted proof built on it is retained verbatim in the delivered body.\par\medskip
\ifdefined\hypertarget\hypertarget{sourcepointer-9}{}\fi
\noindent\textit{Source reference.} Original lemma 0G9H: sources/stacks-chow/tags.txt, source lines 20302--20302. Prerequisite pointer invoked at proof line(s) 1701; it names a standard fact of a sibling chapter that the counted argument uses, and every step of the counted proof built on it is retained verbatim in the delivered body.\par\medskip
\ifdefined\hypertarget\hypertarget{sourcepointer-10}{}\fi
\noindent\textit{Source reference.} Original lemma 0ALR: sources/stacks-chow/tags.txt, source lines 13379--13379. Prerequisite pointer invoked at proof line(s) 1729; it names a standard fact of a sibling chapter that the counted argument uses, and every step of the counted proof built on it is retained verbatim in the delivered body.\par\medskip
\ifdefined\hypertarget\hypertarget{sourcepointer-11}{}\fi
\noindent\textit{Source reference.} Original lemma 0G9G: sources/stacks-chow/tags.txt, source lines 20301--20301. Prerequisite pointer invoked at proof line(s) 1702; it names a standard fact of a sibling chapter that the counted argument uses, and every step of the counted proof built on it is retained verbatim in the delivered body.\par\medskip
\begin{lemma}
\label{lemma-approximate-presentation-rig-smooth}
Let $I$ be an ideal of a Noetherian ring $A$. Let $B$ be an object
of (\ref{equation-C-prime}). Let $B = A[x_1, \ldots, x_r]^\wedge/J$
be a presentation. Assume there exists an element
$b \in B$, $0 \leq m \leq r$, and $f_1, \ldots, f_m \in J$
such that
\begin{enumerate}
\item $V(b) \subset V(IB)$ in $\Spec(B)$,
\item the image of
$\Delta = \det_{1 \leq i, j \leq m}(\partial f_j/\partial x_i)$
in $B$ divides $b$, and
\item $b J \subset (f_1, \ldots, f_m) + J^2$.
\end{enumerate}
Then there exists a finite type $A$-algebra $C$ and an $A$-algebra
isomorphism $B \cong C^\wedge$.
\end{lemma}

\begin{proof}
The conditions imply that $B$ is rig-smooth over $(A, I)$, see
Lemma \ref{lemma-equivalent-with-artin-smooth}.
Write $b' \Delta = b$ in $B$ for some $b' \in B$.
Say $I = (a_1, \ldots, a_t)$. Since $V(b) \subset V(IB)$ there
exists an integer $c \geq 0$ such that $I^cB \subset bB$.
Write $bb_i = a_i^c$ in $B$ for some $b_i \in B$.

\medskip\noindent
Choose an integer $n \gg 0$ (we will see later how large).
Choose polynomials $f'_1, \ldots, f'_m \in A[x_1, \ldots, x_r]$
such that $f_i - f'_i \in I^nA[x_1, \ldots, x_r]^\wedge$.
We set $\Delta' = \det_{1 \leq i, j \leq m}(\partial f'_j/\partial x_i)$
and we consider the finite type $A$-algebra
$$
C = A[x_1, \ldots, x_r, z_1, \ldots, z_t]/
(f'_1, \ldots, f'_m,
z_1\Delta' - a_1^c, \ldots, z_t\Delta' - a_t^c)
$$
We will apply Lemma \ref{lemma-presentation-rig-smooth} to $C$.
We compute
$$
\det\left(
\begin{matrix}
\text{matrix of partials of} \\
f'_1, \ldots, f'_m, z_1\Delta' - a_1^c, \ldots, z_t\Delta' - a_t^c \\
\text{with respect to the variables} \\
x_1, \ldots, x_m, z_1, \ldots, z_t
\end{matrix}
\right) =
(\Delta')^{t + 1}
$$
Hence we see that $\Ext^1_C(\NL_{C/A}, N)$ is annihilated by
$(\Delta')^{t + 1}$ for all $C$-modules $N$. Since $a_i^c$ is
divisible by $\Delta'$ in $C$ we see that $a_i^{(t + 1)c}$ annihilates
these $\Ext^1$'s also. Thus $I^{c_1}$ annihilates
$\Ext^1_C(\NL_{C/A}, N)$ for all $C$-modules $N$
where $c_1 = 1 + t((t + 1)c - 1)$. The exact value of $c_1$ doesn't matter for
the rest of the argument; what matters is that it is independent of $n$.

\medskip\noindent
Since $\NL_{C^\wedge/A}^\wedge = \NL_{C/A} \otimes_C C^\wedge$ by
Lemma \ref{lemma-NL-is-completion} we conclude that multiplication
by $I^{c_1}$ is zero on $\Ext^1_{C^\wedge}(\NL_{C^\wedge/A}^\wedge, N)$
for any $C^\wedge$-module $N$ as well, see
More on Algebra, Lemmas
\ref{more-algebra-lemma-base-change-property-ext-1-annihilated} and
\ref{more-algebra-lemma-two-term-base-change}.
In particular $C^\wedge$ is rig-smooth over $(A, I)$.

\medskip\noindent
Observe that we have a surjective $A$-algebra homomorphism
$$
\psi_n : C \longrightarrow B/I^nB
$$
sending the class of $x_i$ to the class of $x_i$ and sending the class of
$z_i$ to the class of $b_ib'$. This works because of our choices of $b'$
and $b_i$ in the first paragraph of the proof.

\medskip\noindent
Let $d = d(\text{Gr}_I(B))$ and $q_0 = q(\text{Gr}_I(B))$ be the integers
found in 
Local Cohomology, Section \ref{local-cohomology-section-uniform}.
By Lemma \ref{lemma-get-morphism-general-better} if we take
$n \geq \max(q_0 + (d + 1)c_1, 2(d + 1)c_1 + 1)$ we can find a homomorphism
$\varphi : C^\wedge \to B$ of $A$-algebras which is congruent to
$\psi_n$ modulo $I^{n - (d + 1)c_1}B$.

\medskip\noindent
Since $\varphi : C^\wedge \to B$ is surjective modulo $I$
we see that it is surjective (for example use
Algebra, Lemma \ref{algebra-lemma-completion-generalities}).
To finish the proof it suffices to show that
$\Ker(\varphi)/\Ker(\varphi)^2$ is annihilated by a power of $I$, see
More on Algebra, Lemma \ref{more-algebra-lemma-quotient-by-idempotent}.

\medskip\noindent
Since $\varphi$ is surjective we see that
$\NL_{B/C^{\wedge}}^\wedge$ has cohomology modules
$H^0(\NL_{B/C^{\wedge}}^\wedge) = 0$ and
$H^{-1}(\NL_{B/C^{\wedge}}^\wedge) = \Ker(\varphi)/\Ker(\varphi)^2$.
We have an exact sequence
$$
H^{-1}(\NL_{C^\wedge/A}^\wedge \otimes_{C^\wedge} B) \to
H^{-1}(\NL_{B/A}^\wedge) \to
H^{-1}(\NL_{B/C^{\wedge}}^\wedge) \to
H^0(\NL_{C^\wedge/A}^\wedge \otimes_{C^\wedge} B) \to
H^0(\NL_{B/A}^\wedge) \to 0
$$
by Lemma \ref{lemma-exact-sequence-NL}. The first two modules are
annihilated by a power of $I$ as $B$ and $C^\wedge$ are rig-smooth over
$(A, I)$. Hence it suffices
to show that the kernel of the surjective map
$H^0(\NL_{C^\wedge/A}^\wedge \otimes_{C^\wedge} B) \to
H^0(\NL_{B/A}^\wedge)$ is annihilated by a power of $I$.
For this it suffices to show that it is annihilated by
a power of $b$. In other words, it suffices to show that
$$
H^0(\NL_{C^\wedge/A}^\wedge) \otimes_{C^\wedge} B[1/b]
\longrightarrow
H^0(\NL_{B/A}^\wedge) \otimes_B B[1/b]
$$
is an isomorphism. However, both are free $B[1/b]$ modules
of rank $r - m$ with basis $\text{d}x_{m + 1}, \ldots, \text{d}x_r$
and we conclude the proof.
\end{proof}

\end{document}
